% !TEX root = ../main.tex
%% Sujet type bac n° 6 — sujet inédit OptiTAB (valeurs vérifiées : verifications/sujets_bac.py)

\sujetBac{6}{{Le paracétamol : synthèse, structure, dosage spectrophotométrique}/9,{L'énergie d'un satellite : photons, cellules photovoltaïques}/6,{Le niveau sonore d'un concert : intensité sonore, atténuation}/5}

\begin{exoBac}{9}{Le paracétamol}
Le paracétamol est le médicament le plus vendu en France. On étudie sa synthèse au laboratoire, puis le contrôle de la
masse de principe actif d'un comprimé.
\begin{center}
\setchemfig{atom sep=1.9em}
\chemfig{HO-*6(-=-(-NH_2)=-=)}\qquad\qquad
\chemfig{HO-*6(-=-(-N(-[:90]H)-[:-30](=[:-90]O)-[:30])=-=)}\par\smallskip
{\legende{4-aminophénol (à gauche) et paracétamol (à droite).}}
\end{center}

\partieExo{1}{La synthèse}
\begin{docu}{1}{Le protocole}
Dans un erlenmeyer, on introduit \qty{3,0}{g} de 4-aminophénol, \qty{30}{mL} d'eau et \qty{4,0}{mL} d'anhydride éthanoïque, puis on chauffe à
\qty{80}{\degreeCelsius} pendant 20 minutes. Après refroidissement dans un bain d'eau glacée, on filtre sous pression réduite, puis on
recristallise le solide dans l'eau. On obtient \qty{3,3}{g} de cristaux, dont la température de fusion, mesurée au banc
Kofler, vaut \qty{169}{\degreeCelsius}. Équation de la réaction :
\[\ce{C6H7NO + C4H6O3 -> C8H9NO2 + C2H4O2}\]
\end{docu}
\begin{docu}{2}{Données}
{\small\begin{tabular}{@{}lccc@{}}
Espèce & 4-aminophénol & anhydride éthanoïque & paracétamol \\
Masse molaire (\unit{g.mol^{-1}}) & 109,1 & 102,1 & 151,2 \\
Masse volumique (\unit{g.mL^{-1}}) & solide & 1,08 & solide \\
Température de fusion (\unit{\degreeCelsius}) & 187 & -- & 169 \\
\end{tabular}}\par\smallskip
Électronégativités : H 2,2 ; C 2,6 ; N 3,0 ; O 3,4. Économie d'atomes : rapport de la masse molaire du produit recherché à
la somme des masses molaires des réactifs.
\end{docu}
\begin{questions}
\item Recopier la formule du paracétamol, entourer ses groupes caractéristiques et nommer les familles correspondantes.
\item Déterminer le réactif limitant.
\item Calculer le rendement de la synthèse.
\item Calculer l'économie d'atomes de la synthèse. Comment l'améliorer du point de vue de la chimie verte ?
\item Quel est le but de la recristallisation ? Que peut-on conclure de la mesure de la température de fusion ?
\item Lors de la première étape du mécanisme, l'atome d'azote du 4-aminophénol se lie à un atome de carbone d'un groupe
C$=$O de l'anhydride. Identifier le site donneur et le site accepteur de doublet d'électrons, en justifiant.
\end{questions}

\partieExo{2}{Le contrôle d'un comprimé}
\begin{docu}{3}{Le dosage par spectrophotométrie}
Un comprimé « paracétamol 500~mg » est dissous dans une fiole jaugée de \qty{1,000}{L} ; après filtration, la solution est
diluée 50 fois. L'absorbance de la solution diluée, mesurée à \qty{243}{nm} dans une cuve de \qty{1,00}{cm}, vaut $A=0{,}640$. Le
paracétamol est la seule espèce qui absorbe à cette longueur d'onde. L'incertitude-type relative sur la masse déterminée
est estimée à 1,5~\%.
\end{docu}
\schema{bac-s6-etalonnage}
\begin{questions}[resume]
\item Pourquoi effectue-t-on les mesures à \qty{243}{nm}, longueur d'onde du maximum d'absorption du paracétamol ?
\item Énoncer la loi de Beer-Lambert. Les mesures la vérifient-elles ? Déterminer le coefficient de proportionnalité.
\item \initiative Le comprimé contient-il bien \qty{500}{mg} de paracétamol ?
\item Pourquoi a-t-il fallu diluer la solution du comprimé avant la mesure ?
\end{questions}
\end{exoBac}

\begin{exoBac}{6}{L'énergie d'un satellite}
Un satellite en orbite basse tourne autour de la Terre en \qty{98,5}{min}, dont \qty{35}{min} dans l'ombre de la Terre. Ses panneaux solaires,
d'une surface totale de \qty{20}{m^2}, reçoivent en plein soleil une puissance surfacique de \qty{1361}{W.m^{-2}} et convertissent 30~\% de
l'énergie reçue en énergie électrique. Les équipements du satellite consomment en permanence \qty{5,0}{kW} ; pendant
l'éclipse, des batteries prennent le relais.
\begin{donnees}
$h=\qty{6,63e-34}{J.s}$ ; $c=\qty{3,00e8}{m.s^{-1}}$ ; $\qty{1}{eV}=\qty{1,60e-19}{J}$. Dans une cellule au silicium, un photon n'est absorbé que si
son énergie dépasse \qty{1,12}{eV} ; l'excédent d'énergie est perdu sous forme thermique.
\end{donnees}
\begin{questions}
\item Décrire l'effet photoélectrique. Pourquoi a-t-il fallu introduire le modèle du photon pour l'expliquer ?
\item Calculer l'énergie, en joules et en électronvolts, d'un photon de longueur d'onde \qty{550}{nm}.
\item Calculer la longueur d'onde maximale des photons absorbés par le silicium. À quel domaine appartient-elle ?
\item Quelle part de l'énergie d'un photon de \qty{550}{nm} est perdue sous forme thermique ?
\item En supposant que tous les photons reçus aient une longueur d'onde de \qty{550}{nm}, calculer le nombre de photons qui
frappent chaque mètre carré de panneau par seconde.
\item \initiative Les panneaux produisent-ils assez d'énergie pour faire fonctionner le satellite en permanence ?
\end{questions}
\end{exoBac}

\begin{exoBac}{5}{Le niveau sonore d'un concert}
Lors d'un concert en plein air, un sonomètre placé à \qty{10}{m} d'une enceinte mesure un niveau d'intensité sonore de \qty{105}{dB}.
\begin{donnees}
Intensité sonore de référence : $I_0=\qty{1,0e-12}{W.m^{-2}}$ ; $L=10\log\dfrac{I}{I_0}$. Une source ponctuelle de puissance sonore $P$
produit, à la distance $d$ et sans absorption, l'intensité $I=\dfrac{P}{4\pi d^2}$. Réglementation : une exposition à \qty{85}{dB}
est tolérée 8 heures par jour ; chaque augmentation de \qty{3}{dB} divise cette durée par deux.
\end{donnees}
\begin{questions}
\item Calculer l'intensité sonore à \qty{10}{m} de l'enceinte.
\item En assimilant l'enceinte à une source ponctuelle, calculer sa puissance sonore.
\item Montrer que le niveau d'intensité sonore diminue de \qty{6}{dB} chaque fois que la distance double. Calculer le niveau à
\qty{40}{m} de l'enceinte.
\item Deux enceintes identiques, placées côte à côte, fonctionnent ensemble. Calculer le niveau à \qty{40}{m}.
\item \initiative Un spectateur reste deux heures à \qty{40}{m} des deux enceintes. Est-il exposé à un risque ? Des bouchons d'oreille
atténuent le son de \qty{25}{dB} : résolvent-ils le problème ?
\item À quelle distance d'une seule enceinte le niveau descend-il à \qty{85}{dB} ?
\end{questions}
\end{exoBac}
